\section{The selected full-block completion and one-point support comparison}
\label{sec:count-raw-v038}

\status{CRW1--CRW21 reconstruct the actual full-block completed
representation, its relative cotangent presentation, the selected
first-order limit bound, and the one-point support comparison.  The
original limit and patched-module producers and the coefficient and
pro-$\ell$ descent remain explicit inputs to the selected Raw route.
No assertion about arbitrary Kap, vanishing of Sel1, or the complete
Count-II producer is made.}

All formal locators in the accompanying metadata refer to the commit
\par\smallskip\noindent\texttt{7cbf7a0c32fca1badf05931879a37c3f8a2b03b4}.\par\smallskip
The argument below is a human proof on the original objects; the
formal declarations identify its endpoints.  To avoid the existing
class-rank labels CR1--CR6, its claim prefix is CRW.
Write $\mathcal P$ for the source directory
\texttt{AutomorphicVarieties/PatchingModules/}.

\subsection{The actual point, coefficient ring and completed representation}

% BEGIN CRW1
\paragraph{CRW1. The selected scope and its unchanged exponent.}
\label{claim:CRW1}
Fix the original one-point NiceI construction of CCH15, its single
finite coefficient extension and its Kap $k$ with $k.g=0$.
Keep CAP13's realization of the same finite list, its group system
$T$, and its arithmetic realization $\mathrm{TWPl}$.
For Count II put
\[
 n=a=4,\qquad d=[L^+:\mathbb Q],\qquad q=T.q,
 \qquad D=4d,\qquad D_0=D+nq=4d+4q.                 \tag{CRW.1}
\]
CFO20--CFO21 and CCH16 give on these same objects
\[
 \exists C\ \forall m\ \forall^{\mathrm{eventual}}N:
 \quad |\mathrm{Sel1}_{N,m}|
 \leq |\mathrm{Bd1}_{N,m}|\,|K|^{mD_0+C},
 \qquad T.\mathrm{B2}(k).                           \tag{CRW.2}
\]
The finite set assertion in B1 is included.  In B2 the selected
constant is zero and the statement holds at every level and
truncation.  The actual first-order coefficient space at $g=0,a=n$
is $\operatorname{Mat}_n(A)$, where $A=K[[t]]$, and the exact
polarized boundaries are
\[
 dY(\gamma)=\rho(\gamma)Y\rho(\gamma)^{-1}-Y,
 \qquad Y^t=-Y.                                   \tag{CRW.3}
\]
It is the odd coefficient space that is zero, not the first-order
Selmer space.  These identities are CCH1--CCH3 and CFO1.

For the local completion argument allow any original
$h:G\to G_0$, $x:\mathrm{NP}(G_0,n)$ with $a=n>0$, and
$P:\mathrm{OoPt}(h,x,u)$ over an original complete local
Noetherian $O$-algebra $B$.  Here $O$ is the original complete DVR
of residue characteristic $\ell\ne2$, $O\to K$ is surjective,
$f=P.f:B\to A$, and $s=P.s$ satisfies $f(s)=t$.
Put
\[
 \mathfrak p=\ker f,\quad R=\widehat{B_{\mathfrak p}},
 \quad \mathfrak m=\mathfrak m_R,\quad F=K((t)).    \tag{CRW.4}
\]
Every occurrence of $u,P,h,R$ below refers to this one set of
objects.  The original $\mathrm{IsCplAt}$ may be used instead of
the displayed model for the completion: the unique completion
isomorphism carries all its maps and ideals together.
The final application takes $h=\Lambda.h$ for the same supplied
limit $\Lambda:\mathrm{Lim}(T)$ and the same lifted point in the
patching construction.  That supply is not assumed to have been
proved by (CRW.2).
% END CRW1

% BEGIN CRW2
\paragraph{CRW2. Construct the actual Laurent residue and coefficient DVR.}
\label{claim:CRW2}
Choose a uniformizer $\varpi$ of $O$.  Since $s$ cannot be a unit,
$s\in\mathfrak m_B$.  Completeness gives a ring map
\[
 e:O[[S]]\longrightarrow B,\qquad
 \sum b_jS^j\longmapsto\sum b_js^j.
\]
It agrees with the original $O$-map.  Its composite with $f$ is
coefficientwise reduction $O[[S]]\to K[[t]]$.  This does not
presume continuity of an arbitrary ring map: locality of $f$
can be checked directly.
If $b\in\mathfrak m_B$ and $f(b)$ had nonzero constant coefficient,
choose $o\in O^\times$ with the same constant coefficient;
$1-o^{-1}b$ is a unit in $B$, whereas its image has constant
coefficient zero.  This is impossible.  Thus $f$ is local and
the equality of power series follows by taking partial sums.

Coefficientwise lifting proves $f$ is onto and
\[
 B/\mathfrak p=A,\qquad
 \mathfrak m_B=\mathfrak p+(s),\qquad
 \ker(O[[S]]\to A)=(\varpi).                        \tag{CRW.5}
\]
The last equality follows by dividing every coefficient by
$\varpi$.  Consequently $B_{\mathfrak p}$ has residue field $F$,
and completion gives a surjection
\[
 r:R\longrightarrow F,\qquad r(b)=f(b),\quad
 \ker r=\mathfrak m.                               \tag{CRW.6}
\]
This is the original $rF$, not a new residue embedding.

Set
\[
 O^\circ=\widehat{\,O[[S]]_{(\varpi)}\,}.
                                                               \tag{CRW.7}
\]
The localization is a DVR: a nonzero series has a least finite
$\varpi$-order among its coefficients; after division by that
power it is outside $(\varpi)$ and hence is a unit in the
localization.  Its residue field is
$\operatorname{Frac}(K[[t]])=F$.  Noetherianity and the complete
DVR property follow from the checked commutative-algebra imports
listed after CRW21.  The map $e$ carries the complement of
$(\varpi)$ to units of $B_{\mathfrak p}$, and its local map
extends uniquely to
$O^\circ\to R$.  It is local and is an isomorphism on residue
fields.  Thus $R$ is the original $\mathrm{IsCO}\ O^\circ\ R$.

For every $a\in A$, lift its coefficients to $c\in O[[S]]$.
The element $b=e(c)\in B$ and the image $o$ of $c$ in $O^\circ$
satisfy $f(b)=a$ and have the same image in $R$.
Original constants $O$ and the specific section $s=e(S)$ also
come from $O^\circ$.  Fractions of two such lifted series,
with nonzero denominator modulo $\varpi$, show that
$O^\circ\to R\xrightarrow r F$ is onto.  This proves every
coefficient-ring field of the original $\mathrm{Ctx}$:
$\mathrm{oo\_A}$, $\mathrm{oo\_O}$,
$\mathrm{ordinaryPointSection}$ and $\mathrm{oo\_surj}$,
as well as $rF$, $rF\_f$ and $rF\_ker$.
% END CRW2

% BEGIN CRW3
\paragraph{CRW3. A published theorem supplies the actual full determinant representation.}
\label{claim:CRW3}
Let $D_R$ be $u.D$ under the original map $B\to R$.
The residual determinant is exactly
$\det\rho_F$, where
$\rho_F=\rho_x\circ h.p_0$ after $A\to F$.
The source $\mathrm{NP.span}$ says that, after multiplication
by some power of $t$, every matrix in the full even block is
an $A$-linear combination of $\rho_x(G_0')$.
Since $a=n$, that block is all $\operatorname{Mat}_n(A)$.
The surjection $h.p_0$ preserves $G'$; hence
\[
 F\text{-span}\{\rho_F(g):g\in G^0\}
   =\operatorname{Mat}_n(F).                        \tag{CRW.8}
\]
The residual representation is therefore split and absolutely
irreducible: after any field extension its image still spans
the matrix algebra, so a nonzero invariant subspace contains
every vector by applying matrix units.

Form the universal Cayley--Hamilton quotient
$E=R[G^0]/\mathrm{CH}(D_R)$, using all coefficients of the
multivariable Cayley--Hamilton relations, as in the source.
Chenevier, Lemma 1.21, \S1.22 and Theorem 2.22(i), imply
\[
 E\simeq\operatorname{Mat}_n(R),\qquad D_R=\det\widehat\rho
                                                               \tag{CRW.9}
\]
as polynomial laws, including every commutative scalar extension.
All hypotheses have now been checked: $R$ is henselian because
it is complete local; the quotient is Cayley--Hamilton; $n>0$;
and its residual determinant is split absolutely irreducible by
(CRW.8).  This theorem has no invertibility assumption on $n!$
and so applies to $n=4,\ell=3$.

We may and do arrange that $r(\widehat\rho)=\rho_F$ literally.
The supplied residual representation factors through $E\otimes F$,
and is a surjection to $\operatorname{Mat}_n(F)$ by (CRW.8).
It is an isomorphism since (CRW.9) gives dimension $n^2$.
The discrepancy between the two matrix identifications is an
inner automorphism: for matrix units $E_{ij}$, choose nonzero
$v_1$ in the image of the first idempotent and put
$v_i=\phi(E_{i1})v_1$; these vectors are a basis and conjugate
the matrix units.  Lift its change-of-basis matrix to
$\operatorname{GL}_n(R)$, which is possible since $r$ is onto,
and conjugate (CRW.9).  This preserves the full determinant law.
% END CRW3

% BEGIN CRW4
\paragraph{CRW4. Normalize the original polarization over this same ring.}
\label{claim:CRW4}
The source involution of the group algebra sends
\[
 g\longmapsto e(g)^{1-n}c(g)^{-1},
 \qquad e=G.\mathrm{frame.eps}(O,R).
\]
Its square is the identity, by the original conjugation and
cyclotomic identities.  The first coefficient of $\mathrm{Pair.P}$
says that its trace is unchanged.  It consequently descends
through (CRW.9).  Indeed, if $\widehat\rho(z)=0$, then for every
group-algebra element $y$,
\[
 \operatorname{tr}\widehat\rho(y\,z^*)
 =\operatorname{tr}\widehat\rho((y z^*)^*)
 =\operatorname{tr}\widehat\rho(z y^*)=0.
\]
The matrix trace pairing is perfect over every commutative ring:
testing against $E_{ji}$ recovers the $(i,j)$ entry.  Surjectivity
in (CRW.9) therefore makes $\widehat\rho(z^*)=0$.
The resulting anti-involution reduces to transpose, by the
literal form-one polarization of $x$.

An $R$-linear automorphism $\phi$ of $\operatorname{Mat}_n(R)$
reducing to the identity is inner by an elementary construction:
take $v_1=\phi(E_{11})e_1$ and $v_i=\phi(E_{i1})v_1$.
The matrix $V$ with these columns is the identity modulo
$\mathfrak m$, and $\phi(E_{ij})V=VE_{ij}$.
Apply this to the anti-involution followed by transpose.
It has the form $X\mapsto H^{-1}X^tH$ with $H\equiv1$.
Squaring it shows $H^t=cH$ for a central unit $c$ with
$c^2=1$ and $c\equiv1$.  Since $2$ is a unit, $c+1$ is a unit
and $c=1$.  Thus $H$ is symmetric.

There is a symmetric $Q\equiv1$ commuting with $H$ such that
$Q^tHQ=1$.  One explicit construction starts with $U_0=1$ and
iterates $U_{j+1}=(U_j+H U_j^{-1})/2$.
All $U_j$ are symmetric units commuting with $H$, and their
squared errors $U_j^2-H$ double their $\mathfrak m$-adic order.
Completeness gives $U^2=H$; set $Q=U^{-1}$.
Conjugating the representation by $Q$ gives
\[
 \widehat\rho(g)^t\widehat\rho(cg)=e(g)^{1-n}1,
 \qquad r(\widehat\rho(g))=\rho_F(g).                \tag{CRW.10}
\]
No factorial or square-root extension of the coefficient ring
has been used.  This is the original form, with exponent $1-n$.
% END CRW4

% BEGIN CRW5
\paragraph{CRW5. A finite original-$B$ lattice and uniform continuity.}
\label{claim:CRW5}
Choose $n^2$ group elements whose residual matrices form an
$F$-basis of $\operatorname{Mat}_n(F)$.  Their lifts $E_i$ form
an $R$-basis because their determinant is a unit.  The trace
Gram matrix $\Delta_{ij}=u.D.\mathrm{tr}(g_i g_j)$ has entries
in $B$ and determinant $\delta\notin\mathfrak p$.
For every $g$, the coordinates of $\widehat\rho(g)$ are
\[
 \Delta^{-1}
 \big(u.D.\mathrm{tr}(g g_j)\big)_j.
\]
They belong to the image of $\delta^{-1}B$ in $R$.
Thus all group matrices lie in the finite $B$-module
$P_0=\sum B\delta^{-1}E_i$.
The module $M$ spanned over $B$ by the group matrices is
finite, since $B$ is Noetherian and $M\subset P_0$.
It contains $1$ and is a $B$-algebra, since products of group
matrices are again group matrices.

For every $j$ the original continuity of $u.D$ and compactness
of $G^0$ give an open normal $V$ such that all differences
$u.D.\mathrm{tr}(gv g_i)-u.D.\mathrm{tr}(g g_i)$ lie in
$\mathfrak m_B^j$, uniformly for $g\in G^0,v\in V$.
To see the uniformity, reduce each continuous trace function
modulo $\mathfrak m_B^j$.  A continuous function from a compact
profinite group to a discrete set has finite image and is
constant on cosets of one open normal subgroup; intersect
these finitely many subgroups and their conjugates as needed.
The coordinate formula gives the corresponding matrix
differences in $\mathfrak m_B^jP_0$.

For the finite inclusion $M\subset P_0$, Artin--Rees gives
$c$ with
$M\cap\mathfrak m_B^{j+c}P_0\subset\mathfrak m_B^jM$.
Apply the previous argument at $j+c$.  We obtain the exact
source continuity field
\[
 \forall j\ \exists V\text{ open normal}\ \forall g\ \forall v\in V,
 \quad\widehat\rho(gv)-\widehat\rho(g)\in\mathfrak m_B^jM.
                                                               \tag{CRW.11}
\]
Here $M$ is one fixed finite module for all $j$, not a different
lattice at each precision.  It also satisfies
$\bigcap_j\mathfrak m_B^jM=0$ by Krull intersection, whose
hypotheses are the local Noetherian ring $B$ and this finite
module.
% END CRW5

% BEGIN CRW6
\paragraph{CRW6. The whole original closed ideal and every URep field.}
\label{claim:CRW6}
Let $Z$ be the exact source set
$\mathrm{pairGens}(u)\cup\mathrm{qGens}(P)$, including its
star images, and let
\[
 I=P.\mathrm{ker}=\mathrm{closureIdeal}_{\mathfrak m_B}
                  (\mathrm{CH}(u.D)\cup Z).
\]
Every generator in $Z$ acts by zero under $\widehat\rho$.
For the ordinary products and the inertia differences, this
follows from $\mathrm{Pair.Ord}$ and $\mathrm{Pair.I}$:
their products with every group element have trace zero,
so CRW3 and the perfect matrix trace pairing make the matrix
itself zero.  The same argument uses $P.q1$ and $P.q2$ for
the two auxiliary-prime relations.  Their star images vanish
by CRW4.  Every multivariable Cayley--Hamilton coefficient
vanishes by (CRW.9) and the ordinary matrix identity over
the polynomial scalar ring.

The word ``closed'' here retains both indices in the original
definition.  For each integer $j$ and open normal subgroup
$V\subset G^0$, take the two-sided ideal generated by these
relations, by coefficient elements in $\mathfrak m_B^j$,
and by $v-1$ for $v\in V$; $I$ is their intersection.
Choose $V$ as in (CRW.11).  Images of its generators lie in
$\mathfrak m_B^jM$, and this is a two-sided ideal of the
finite algebra $M$.  Therefore the image of $I$ lies in
every $\mathfrak m_B^jM$ and is zero.  This proves the literal
$\mathrm{kills}$ field, including the group-neighborhood part
of the closure.

There is also an equality of the actual quotients
\[
 R[G^0]/\langle\operatorname{image} I\rangle
       \simeq\operatorname{Mat}_n(R).               \tag{CRW.12}
\]
Indeed $I$ contains all original Cayley--Hamilton coefficients.
Their ideal after scalar extension is $\mathrm{CH}(D_R)$:
the universal multivariable coefficient construction commutes
with arbitrary base change, as in Chenevier \S1.22.
This gives one inclusion between the kernels; the vanishing
just proved gives the reverse inclusion through (CRW.9).
For any field map $R\to\kappa$, taking the quotient and
extending scalars consequently gives
$\kappa[G^0]/\langle\operatorname{image}I\rangle
\simeq\operatorname{Mat}_n(\kappa)$.
This ring is semisimple: its left regular module is the direct
sum of its $n$ column modules, each simple by matrix units.
Thus the original $\mathrm{NonSS}(I,B\to R,Q)$ is false for
every prime $Q$ of $R$, including every field extension
allowed by that definition.

Take $AA=R$, the identity structure map, the zero odd module
$N=0$, representation (CRW.10), and lattice $M$ of CRW5.
It is finite over $R$, its scalar map is injective,
$R\cdot1\oplus0=R$, and the odd-product condition is vacuous.
Its span is all matrices by CRW3, which is exactly
$\mathrm{blkMat}(n,n,1,0)$.
The determinant, form, reduction, kills and continuity fields
have been proved.  Its $\mathrm{thick}$ field holds because
NonSS is impossible; it is not deduced from the false
inclusion $\operatorname{Ann}(0)\subset Q$.
This constructs the full original $\mathrm{URep}$, hence the
actual $a=n$ instance of S6.  It uses no general reducible
representation or nonzero odd-part producer.
% END CRW6

\subsection{The original cotangent space injects into first-order Selmer classes}

% BEGIN CRW7
\paragraph{CRW7. Relative cotangent duality on the actual completion.}
\label{claim:CRW7}
Use the coefficient DVR $O^\circ$ from CRW2 and put
\[
 J_2=\mathfrak m_{O^\circ}R+\mathfrak m^2,
 \qquad V=\mathfrak m/J_2.
\]
This is a finite-dimensional $F$-space: $\mathfrak m$ is
finitely generated and multiplication by $\mathfrak m$ is zero
on the quotient.  The field action is the one supplied by
$r:R\to F$.  Let $\mathcal T$ be the $F$-space of additive
maps $\alpha:R\to F$ satisfying
\[
 \alpha(bc)=r(b)\alpha(c)+r(c)\alpha(b),
 \qquad \alpha(O^\circ)=0.                         \tag{CRW.13}
\]
Restriction gives an isomorphism
\[
 \mathcal T\simeq\operatorname{Hom}_F(V,F).          \tag{CRW.14}
\]
Here is the complete inverse.  Given $\lambda:V\to F$ and
$b\in R$, choose $o\in O^\circ$ with $r(o)=r(b)$, and set
$\alpha(b)=\lambda([b-o])$.
A second choice differs by the image of
$\mathfrak m_{O^\circ}$, so the value is independent of choice.
Addition is immediate.  Writing $b=o+x,c=o'+y$ with
$x,y\in\mathfrak m$, their product differs from $oo'$ by
$oy+o'x+xy$; the last term is zero in $V$.
This proves the product rule.  Conversely a derivation
annihilates $\mathfrak m^2$ and
$\mathfrak m_{O^\circ}R$ by (CRW.13), and its restriction is
$F$-linear.  The two constructions are inverse.

For the URep of CRW6, $N=0$ implies
$\mathrm{ired}(N)=0$.  Thus $J_2$ and $\mathcal T$ are
exactly the source $\mathrm{Ctx.J2}$ and $\mathrm{Ctx.Tp}$;
no larger quotient or smaller tangent space has been substituted.
% END CRW7

% BEGIN CRW8
\paragraph{CRW8. Bound the derivative using the original parameter.}
\label{claim:CRW8}
Let $z_1,\ldots,z_r$ generate $\mathfrak p$ as a $B$-ideal.
For $\alpha\in\mathcal T$, set
$L_\alpha=\sum A\alpha(z_i)\subset F$.
For any $b\in B$, choose $c\in O[[S]]$ with $f(e(c))=f(b)$.
Then $b-e(c)=\sum b_i z_i$ and $\alpha(e(c))=0$, since the
same element in $R$ lies in the image of $O^\circ$.  Hence
\[
 \alpha(b)=\sum f(b_i)\alpha(z_i)\in L_\alpha.
\]
In particular this bound is uniform in $b$.  The product rule
and (CRW.5) imply, for $j\ge1$,
\[
 f(\mathfrak m_B^j)\subset t^jA,
 \qquad \alpha(\mathfrak m_B^j)
          \subset t^{j-1}L_\alpha.                 \tag{CRW.15}
\]
One may check the second formula on products of $j$ factors
of $\mathfrak m_B$: in each derivative summand all other
$j-1$ factors contribute a factor $t$.

Choose finite $B$-generators $P_i$ for the matrix lattice $M$
of CRW5.  Entrywise differentiation of
$bP_i$ is $f(b)\alpha(P_i)+\alpha(b)r(P_i)$.
The finitely many matrices $\alpha(P_i)$ and $r(P_i)$ and
(CRW.15) show that $\alpha(M)$ is contained in one finite
$A$-lattice in $\operatorname{Mat}_n(F)$, and that
$\alpha(\mathfrak m_B^jM)$ tends uniformly to zero in the
$t$-adic topology.  Equation (CRW.11) now proves that
$\alpha(\widehat\rho(g))$ is bounded and continuous.
The original $\rho_x$ and its inverse have entries in $A$
and are adically continuous.  Therefore
\[
 z_\alpha(g)=\alpha(\widehat\rho(g))\rho_F(g)^{-1}
                                                               \tag{CRW.16}
\]
is bounded and continuous, with one denominator bound for
all $g$.  This constructs the needed Laurent boundedness on
the original ring; it is not an extra TrBnd assumption.
% END CRW8

% BEGIN CRW9
\paragraph{CRW9. Ordered product relations give the required first-order flag.}
\label{claim:CRW9}
We record the elementary local argument in the precise form
needed here.  Let $F$ be an infinite field,
$E=F[\varepsilon]/(\varepsilon^2)$,
$R_0,R_1:H\to\operatorname{GL}_n(E)$ with the same reduction,
and ordered characters $\chi_{0i},\chi_{1i}$ with the same
reduction.  Suppose $R_0,\chi_0$ are constant over $F$,
$w^{-1}R_0(g)w$ is upper triangular with diagonal
$\chi_{0i}(g)$, these characters are pairwise distinct, and
\[
 \prod_{i=1}^n\big(R_1(g_i)-\chi_{1i}(g_i)1\big)=0
 \quad\hbox{in that order for every }(g_i).         \tag{CRW.17}
\]
There exists a matrix $Y\in\operatorname{Mat}_n(F)$ such that
\[
 w^{-1}\big(z(g)-dY(g)\big)w
 \text{ is upper triangular, with diagonal }
 \frac{\dot\chi_i(g)}{\chi_{0i}(g)},                \tag{CRW.18}
\]
where $R_1=(1+\varepsilon z)R_0$ and
$\chi_{1i}=\chi_{0i}+\varepsilon\dot\chi_i$.

To prove this, extend (CRW.17) to arbitrary group-algebra
arguments by multilinearity in its $n$ arguments.
There is a finite $F$-linear combination $u$ of group elements
for which the values $a_i=\chi_{0i}(u)$ are distinct.
Each equality $a_i=a_j$ defines a proper linear hyperplane
in the span of finitely many elements distinguishing all
pairs; finitely many proper hyperplanes cannot cover a vector
space over the infinite field $F$.
For $b=0,1$ set
\[
 f_{bi}=\prod_{j\ne i}
       \frac{R_b(u)-\chi_{bj}(u)1}
            {\chi_{bi}(u)-\chi_{bj}(u)}.
\]
The denominators are units.  Polynomial division by
$\prod_i(X-\chi_{bi}(u))$ proves that the $f_{bi}$ are
complete orthogonal idempotents, with the same reductions.
Multiply the extended ordered relation, with all arguments
$u$ except the $i$-th, on the left by $f_{ba}$ and on the
right by $f_{bb'}$.  If $a=b'=i$, all outside factors are
units on these idempotents and give
$f_{bi}R_b(g)f_{bi}=\chi_{bi}(g)f_{bi}$.
If $a>b'$, put the variable argument in position $a$:
the left factors have indices below $a$, the right factors
have indices above $a>b'$, so again they are units.  Orthogonality
gives $f_{ba}R_b(g)f_{bb'}=0$.

Put $W=\sum_i f_{1i}f_{0i}$.  It is $1$ modulo
$\varepsilon$, hence a unit, and $Wf_{0i}=f_{1i}W$.
Thus $W^{-1}R_1W$ is upper triangular relative to the
original ordered idempotents, with diagonal $\chi_{1i}$.
Those original idempotents are themselves upper triangular
in the $w$ basis; the vectors $f_{0i}e_i$ form a unit upper
triangular basis.  Consequently the same triangular and
diagonal assertions hold in the original $w$ basis.
Write $W=1+\varepsilon U$ and put $Y=-U$.
Then
$W^{-1}R_1W=(1+\varepsilon(z-dY))R_0$.
Comparing diagonal entries gives (CRW.18), including its sign.
% END CRW9

% BEGIN CRW10
\paragraph{CRW10. Differentiate into the literal polarized Selmer condition.}
\label{claim:CRW10}
The map
$b\mapsto r(b)+\varepsilon\alpha(b)$ is a ring map
$R\to F[\varepsilon]/(\varepsilon^2)$ by (CRW.13).
Apply it to the original $\widehat\rho$.
Multiplication gives the cocycle identity
\[
 z_\alpha(gh)=z_\alpha(g)+
       \rho_F(g)z_\alpha(h)\rho_F(g)^{-1}.
\]
Differentiate (CRW.10).  Its multiplier comes from the
original $O$ and is killed by $\alpha$, so
$z_\alpha(cg)=-z_\alpha(g)^t$.
The original inertia relations in CRW6 give
$\widehat\rho(i)=1$ for every $i\in G.\mathrm{frame.Iu}$;
their derivatives vanish.  This uses exactly
$P.\mathrm{Ip}$ and the source $\mathrm{Pair.I}$.

At each ordinary place the matrix relation in CRW6 is
(CRW.17), with the same ordered characters $u.\psi$.
The source $\mathrm{NP.N1}$ gives their pairwise
distinctness over $F$.  The original flag is $x.w$,
transported through $h.\phi$; no new order is chosen.
CRW9 supplies a local $Y_v$ giving the source
$\mathrm{LocAt}$ condition.
On $\mu_v$, $u.\mathrm{psi.mu}$ has values in the
original $O$-image, so its derivative is zero.
Thus the diagonal in (CRW.18) vanishes on precisely
the subgroup required by $\mathrm{LocAt}$.
Together with CRW8 this proves
\[
 z_\alpha\in\mathrm{SelP}(h,x,F),\qquad
 z_\alpha\text{ is bounded and adically continuous}.\tag{CRW.19}
\]
These are the full source conditions: module membership,
adic continuity, cocycle identity, conjugation, inertia,
and one permitted local correcting matrix at every ordinary
place.  In the selected $g=0,a=n$ case they identify with
the constant-component form of the actual first-order Sel1.
The construction $\alpha\mapsto z_\alpha$ is $F$-linear.
% END CRW10

% BEGIN CRW11
\paragraph{CRW11. The actual generator relations make the cotangent map injective.}
\label{claim:CRW11}
Suppose $z_\alpha=dY$ for a polarized boundary
$Y^t=-Y\in\operatorname{Mat}_n(F)$.
Conjugating the dual-number representation by
$1+\varepsilon Y$ on the left and
$1-\varepsilon Y$ on the right makes it identically
$\rho_F$: indeed
$z_\alpha+Y-\rho_FY\rho_F^{-1}=0$.
The full determinant equality in CRW3 implies that
$\alpha$ kills every $u.D.\lambda_i(g)$.

It also kills the actual ordered characters.
For a fixed index $i$ and local element $g$, use the
ordinary mixed relation with $g$ in position $i$.
For every $j\ne i$ choose an element $g_j$ separating
$\psi_i$ from $\psi_j$, possible by N1.
After the representation has become constant, the $i$-th
diagonal entry of the ordered product is
\[
 (\psi_i(g)-\widetilde\psi_i(g))
 \prod_{j\ne i}
 (\psi_i(g_j)-\widetilde\psi_j(g_j))=0.
\]
All factors in the product are units in the dual-number
ring.  Therefore $\widetilde\psi_i(g)=\psi_i(g)$,
so $\alpha(u.\psi_i(g))=0$.

Fix an auxiliary place $j$.  The distinct eigenvalues
$f(U_{ji})$ and the exact factorization $P.\mathrm{eig}$
give a nonzero $F$-eigenvector for each eigenvalue.
Apply the constant-conjugated q1 relation to that vector.
Every factor except the $i$-th is a unit, and a nonzero
coordinate of the vector is a unit in the dual-number
ring.  Thus $\widetilde U_{ji}=f(U_{ji})$ and
$\alpha(U_{ji})=0$.
The q2 relation on the same vector has all the other
eigenvalue factors invertible and the remaining factor
$1-1-\widetilde y_{ji}$, since $h.p_0(ta_j)=1$.
It gives $\widetilde y_{ji}=0$, hence
$\alpha(y_{ji})=0$.
Finally $\alpha(s)=0$ by CRW2, and $\alpha(O)=0$.

These are exactly all generators in $P.\mathrm{gen}$.
For any $b\in B$ and any $j$, choose an element $b_j$ in
their $O$-algebra with $b-b_j\in\mathfrak m_B^j$.
Then $\alpha(b_j)=0$, and (CRW.15) gives
$\alpha(b)\in t^{j-1}L_\alpha$ for all $j$.
A finite $A$-lattice in $F$ is separated, so $\alpha(b)=0$.
It follows that $\alpha$ vanishes on $B_{\mathfrak p}$,
using the derivative formula for an inverse.
Every element of its completion agrees modulo
$\mathfrak m^2$ with an element of $B_{\mathfrak p}$;
derivations to $F$ kill $\mathfrak m^2$.
Thus $\alpha=0$ on $R$.
We have proved an injection of $F$-spaces
\[
 \mathcal T\lhook\joinrel\longrightarrow
 \frac{\{\text{bounded Laurent SelP cocycles}\}}
      {\{dY:Y^t=-Y\}}.                            \tag{CRW.20}
\]
The proof used the original full polynomial determinant,
ordered characters, both auxiliary relations and the
original parameter; trace equality alone would not justify
the first generator step in residue characteristic three.
% END CRW11

\subsection{The selected first-order bound at the actual limit}

% BEGIN CRW12
\paragraph{CRW12. Finite cocycles stabilize on a supplied original limit.}
\label{claim:CRW12}
Keep the original $T$ and one supplied $\Lambda:\mathrm{Lim}(T)$.
Fix $m\ge1$, write $A_m=A/(t^m)$, and use the constant
component of the selected $g=0$ first-order representation.
There is a set of indices belonging to $\Lambda.U$ on which
pullback to the same group $\Lambda.\mathrm{Ph}^0$ gives
bijections
\[
 \mathrm{Sel1}(\Lambda.h,A_m,k)
     \simeq\mathrm{Sel1}(T.h_N,A_m,k),\qquad
 \mathrm{Bd1}(\Lambda.h,A_m,k)
     \simeq\mathrm{Bd1}(T.h_N,A_m,k).               \tag{CRW.21}
\]
The maps identify cocycles as functions after pullback;
they do not choose a new arithmetic level system.

Here are the finite and uniform details.  Put
$\Pi=\Lambda.\mathrm{Ph}$ and
$\Pi^0=\ker(\eta_G\circ\pi_o)$.
The fields $\mathrm{modH}$ and $T.H_\eta$ show
$\eta_{G_N}\pi_N=\eta_G\pi_o$.
Thus both $\pi_N^0:\Pi^0\to G_N^0$ and
$\pi_o^0:\Pi^0\to G^0$ are continuous surjections.
The subgroup $\Pi^0$ is open of finite index.  It is
topologically finitely generated: take a finite generating
set for $\Pi$ and a finite coset transversal; the finitely
many Schreier words generate the subgroup in the abstract
dense group, and taking closure gives $\Pi^0$.

For a fixed action on a finite matrix module, a continuous
cocycle is determined by its values on those finitely many
generators, by its multiplication formula and continuity.
Hence there are only finitely many such cocycles on $\Pi^0$.
The original $p_{\mathrm{lim}}$ first makes the two pulled-back
actions identical, on one set in $\Lambda.U$: apply it to
an open normal subgroup of $G_0.\Gamma$ on which the
finite representation is trivial.
For each of the finitely many cocycles $z$, the map
$g\mapsto(z(g),\rho(g))$ to the finite semidirect-product
group has an open normal kernel in $\Pi^0$.
Take its normal core in $\Pi$ and then the intersection
over all these cocycles.  This gives one open normal
$V\subset\Pi^0$ on whose cosets all these functions
and the action are constant.

The literal $\mathrm{ker\_lim}$ identity says, eventually
in $\Lambda.U$,
\[
 \ker\pi_N\,V=\ker\pi_o\,V.
\]
For a cocycle with this fixed action, descent through
$\pi_N^0$ means vanishing on $\ker\pi_N$; descent through
$\pi_o^0$ means vanishing on $\ker\pi_o$.
The kernels lie in $\Pi^0$ and act trivially on the module.
Since $z$ is zero on $V$, the displayed equality makes
these two descent conditions equivalent.
Surjectivity makes the descended function unique, and the
quotient topology makes it continuous.

Use the same $V$ in the source's uniform
$\phi_{\mathrm{lim}}$ and $\phi^u_{\mathrm{lim}}$ clauses.
The definition of $\mathrm{Close}$ supplies representatives
with quotient in $V$, so the descended functions have
equal values on every corresponding local element,
simultaneously for all places and all those elements.
The same $ct$ represents complex conjugation in both groups.
Consequently the polarized identity, inertia vanishing and
all local upper triangular and diagonal conditions transfer
in both directions; their witnesses $Y_v$ are the very same
matrices over $A_m$.  The module condition is unchanged.
Adic continuity here is exactly continuity into a finite
discrete module, since $t^m=0$.
This proves the Sel1 bijection.  Equal pulled-back actions
identify the functions $dY$ for every $Y^t=-Y$, proving the
boundary bijection.  Every set intersected in this argument
is finite; no intersection over infinitely many cocycles or
levels has been taken.
% END CRW12

% BEGIN CRW13
\paragraph{CRW13. Transfer B1 with its original constant before all truncations.}
\label{claim:CRW13}
Let $C$ be the single constant in (CRW.2).
For each fixed $m\ge1$, intersect the set in
$\Lambda.U$ furnished by CRW12 with the eventual set
in the definition of B1.  It belongs to $\Lambda.U$
because $\Lambda.U\le\mathrm{atTop}$, and so is nonempty.
Choose any index in it and apply the two bijections.
The result is finiteness and
\[
 |\mathrm{Sel1}(\Lambda.h,A_m,k)|
 \leq |\mathrm{Bd1}(\Lambda.h,A_m,k)|\,
       |K|^{mD_0+C}\qquad\hbox{for every }m.         \tag{CRW.22}
\]
For $m=0$ both coefficient and function spaces have one
element, so the assertion also holds.
The chosen index may depend on $m$, as the original B1
permits.  The exponent constant is still the same $C$,
chosen before $m$.  This is the actual
$\mathrm{FirstSelmerCardBoundAtAllLevels}$ on
$\Lambda.h$; it is not an eventual statement in $m$.
% END CRW13

% BEGIN CRW14
\paragraph{CRW14. Bound approximate boundaries by finitely many DVR equations.}
\label{claim:CRW14}
Let $z_1,\ldots,z_e$ be integral first-order Selmer
cocycles on an arbitrary fixed group $G^0$, with the
selected full coefficient space, and suppose
\[
 \sum c_i z_i=dY, c_i\in A, Y^t=-Y\in\operatorname{Mat}_n(A)
       \quad\Longrightarrow\quad c_1=\cdots=c_e=0.\tag{CRW.23}
\]
There is an integer $b$, independent of $m$, such that
\[
 \sum c_i z_i\pmod {t^m}\in\mathrm{Bd1}(h,A_m,k)
       \quad\Longrightarrow\quad
       c_i\in t^{\max(m-b,0)}A\text{ for every }i.  \tag{CRW.24}
\]

To prove this, use the $n(n-1)/2$ upper off-diagonal
entries as coordinates on skew matrices.  This is exact
because $2$ is a unit; their diagonal entries vanish.
The equations
$\sum c_i z_i(g)-dY(g)=0$, for all $g$ and all entries,
are linear equations over $A$ in the finite list
$(c_i,Y_{ab})$.  Their rows generate a submodule of a
finite free $A$-module.  Since a DVR is Noetherian,
finitely many of those rows generate it.  Let $E$ be
the finite matrix formed by such rows.
Vanishing of $E$ is equivalent to all the original
equations, over $A$, $F$, and every $A_m$.

Over a DVR elementary invertible row and column
operations diagonalize $E$, with diagonal entries
$t^{b_1},\ldots,t^{b_r}$ and zero entries elsewhere.
Unit factors have been absorbed.  One can obtain this by
choosing an entry of least valuation, using it to
clear its row and column, and inducting on the
remaining size.  Let $b=\max(0,b_1,\ldots,b_r)$.
For a vector $v$ with $Ev\in t^mA^r$, its coordinates
in the nonzero diagonal directions lie in
$t^{\max(m-b,0)}A$.  The remaining component lies in
$\ker E$.  Every vector in that kernel has all
$c$-coordinates zero by (CRW.23).
The invertible column operations have entries in $A$,
so the $c$-coordinates of $v$ satisfy the same bound.

Finally a skew matrix over $A_m$ lifts to a skew
matrix over $A$ by lifting its upper off-diagonal
entries and taking negatives below the diagonal.
Apply the preceding bound to this lift of the
boundary witness.  This proves (CRW.24) without
assuming that a modular boundary lifts as an exact
integral boundary.
% END CRW14

% BEGIN CRW15
\paragraph{CRW15. The exact integral K1 and its Laurent version.}
\label{claim:CRW15}
Use the group $G=\Lambda.G$ and (CRW.22).
For independent integral classes as in (CRW.23),
the map
\[
 (A/t^m)^e\longrightarrow
 \mathrm{Sel1}(\Lambda.h,A_m,k)\big/
                         \mathrm{Bd1}(\Lambda.h,A_m,k)
\]
has kernel of size at most $|K|^{eb}$ by CRW14,
for $m\ge b$.  Boundaries are contained in Sel1:
the cocycle identity is direct, the skew condition
gives the polarized identity, inertia acts trivially,
and the local correcting matrix can be the same $Y$.
All these spaces are modules, and (CRW.22) makes
their finite quotient meaningful.  Therefore
\[
 |K|^{e(m-b)}
 \leq \frac{|\mathrm{Sel1}(\Lambda.h,A_m,k)|}
               {|\mathrm{Bd1}(\Lambda.h,A_m,k)|}
 \leq |K|^{mD_0+C}.
\]
Since $|K|>1$, division by $m$ and arbitrarily large
$m$ give $e\le D_0$.  This is the original K1,
including its independence modulo actual boundaries.

More generally, on any fixed original $h$ for which
K1 holds, the same bound holds for any finite family
of bounded Laurent SelP cocycles independent modulo
the polarized Laurent boundaries.  Multiply the whole family by a
common power of $t$ which makes their values integral
and makes all their local correcting matrices integral.
Such a power exists by boundedness and because the
set of ordinary places is finite.  Adic continuity
and every defining equation are preserved.
These cocycles give the literal selected Sel1 over
$A$ via its constant component.  An integral
linear relation which is an integral boundary would
also be a Laurent relation, so its coefficients vanish.
Applying K1 gives the bound.  This proves precisely
the full-block instance of $\mathrm{k1F}$.

On this same selected Kap, K2 and K3 are also supplied.
For K2 the only coefficient index is the singleton
$\operatorname{Fin}(g+1)=\operatorname{Fin}(1)$;
take its coefficient to be $1$ and take $\Psi=0$.
All equations indexed by $\operatorname{Fin}(g)$ are
empty.  For K3 the literal odd matrix module is zero
when $a=n$, so its only cocycle is zero; take $s=0$
and the empty coefficient family.  Neither argument
says that Sel1 is zero.
Thus the selected CohLim endpoint is proved from
the same supplied $\Lambda$ and B1, with the chosen
Kap fixed throughout.
% END CRW15

\subsection{Actual generators, the dimension comparison and one-point support}

% BEGIN CRW16
\paragraph{CRW16. Construct the original cotangent generators.}
\label{claim:CRW16}
Apply CRW7--CRW11 to the actual $P,R,O^\circ$ of CRW1,
and assume K1 for its same map $h$ and selected $k$.
Let $e=\dim_F\mathcal T$ and choose a basis
$\alpha_1,\ldots,\alpha_e$.
If a linear combination of the corresponding $z_{\alpha_i}$
were a polarized Laurent boundary, CRW11 would make
$\sum c_i\alpha_i=0$, so every coefficient would be zero.
The Laurent conclusion of CRW15 gives $e\le D_0$.
By (CRW.14), $\dim_FV=e$ as well.

Lift an $F$-basis of $V$ to elements
$X_1,\ldots,X_e\in\mathfrak m$, and append zero
elements to make exactly $D_0$ elements.  Lifting
the scalar coefficients through $r$ shows
\[
 \mathfrak m\subset (X_1,\ldots,X_{D_0})
               +\mathfrak m_{O^\circ}R+\mathfrak m^2.
                                                               \tag{CRW.25}
\]
For the actual full-block URep, $N=0$ and
$\mathrm{ired}(N)=0$ by CRW6--CRW7.
The odd generators indexed by $\operatorname{Fin}(g)$
and all their quadratic coefficient arrays are empty
for the chosen $g=0$.
Consequently (CRW.25), together with those empty
families, proves every field of the original
$\mathrm{Ctx.Pres}(D_0)$ in this scope: maximal-ideal
membership of $X$, generation of $N$, the independence
of the empty quadratic family, the displayed
cotangent generation and all empty odd relations.
This is an actual construction of the source
$\mathrm{pres}$ output, not the presentation inequality
added as a new hypothesis.
% END CRW16

% BEGIN CRW17
\paragraph{CRW17. The actual complete ring is a power-series quotient.}
\label{claim:CRW17}
This paragraph proves the commutative-algebra step
with precisely the original coefficient ring.
Let $\pi^\circ$ be a uniformizer of $O^\circ$ and
$I=(\pi^\circ,X_1,\ldots,X_{D_0})R$.
Equation (CRW.25) says $\mathfrak m=I+\mathfrak m^2$.
The finite module $\mathfrak m/I$ thus equals
$\mathfrak m(\mathfrak m/I)$ and is zero by Nakayama.
One elementary proof of the latter assertion writes a
finite generating column $v$ as $Cv$ with entries of
$C$ in $\mathfrak m$; the adjugate of $1-C$ and its
unit determinant give $v=0$.
We have proved $\mathfrak m=I$.

Evaluation in the $X_i$ gives a continuous map
\[
 \Phi:O^\circ[[Z_1,\ldots,Z_{D_0}]]\longrightarrow R,
 \qquad Z_i\longmapsto X_i.                        \tag{CRW.26}
\]
It is onto.  First lift the residue of any $b\in R$
through $O^\circ\to F$.  If an approximation has
error in $\mathfrak m^j$, express that error modulo
$\mathfrak m^{j+1}$ by degree-$j$ monomials in
$\pi^\circ,X_i$ with coefficients lifted from $F$.
Lift this homogeneous correction to degree $j$ in
the maximal ideal $(\pi^\circ,Z_i)$ of the source.
The successive corrections converge there: for each
monomial in the $Z_i$, the coefficients form a
$\pi^\circ$-adically convergent series, and the tails
have arbitrarily high total maximal-ideal order.
Completeness of $R$ makes their image equal $b$.
This proves surjectivity without assuming an algebraic
section of the Laurent residue field.
% END CRW17

% BEGIN CRW18
\paragraph{CRW18. The full-block local prime comparison, with the exact dimension.}
\label{claim:CRW18}
The source of (CRW.26) is a Noetherian local domain
of dimension $1+D_0$.
For the dimension upper bound its maximal ideal has
$1+D_0$ generators, so the checked local dimension
theorem gives at most that number.
For the lower bound use the strict prime chain
\[
 (0)\subset(\pi^\circ)\subset
 (\pi^\circ,Z_1)\subset\cdots\subset
 (\pi^\circ,Z_1,\ldots,Z_{D_0}).                  \tag{CRW.27}
\]
Primality follows from the remaining power-series
domain after each quotient.  Power series over a
domain form a domain, by taking the first nonzero
coefficient in a multiplicative monomial order.

Let $Z$ be any prime of $R$ for which
$\dim(R/Z)\ge1+D_0$.  If $\Phi^{-1}(Z)$ were
nonzero, any chain of primes of length $1+D_0$ in
the quotient would lift to such a chain starting
at a nonzero prime of the source; prepending its
zero prime contradicts (CRW.27) and the upper
dimension bound.  Thus $\Phi^{-1}(Z)=0$.
Surjectivity of $\Phi$ gives $Z=0$, and also
$\ker\Phi=0$.  In particular
\[
 \dim(R/Z)\ge1+D_0\quad\Longrightarrow\quad
 Z=0\subset Q\quad\hbox{for every prime }Q.         \tag{CRW.28}
\]
This is the complete proof of the actual
$\mathrm{presOut0}$ application.  It even identifies
$R$ with the regular power-series ring whenever
such a $Z$ exists; that existence is not asserted
by the cotangent argument.

CRW2--CRW17 therefore give the $a=n$, selected-$g=0$
instance of $\mathrm{GalOut}$ on the original
completion.  It needs K1 for that same $h,k,D_0$.
The source's K2, K3, stable-discriminant and target
stable-or-thick antecedents can be retained, but
(CRW.28) does not use them.  The required dimension
for Count II remains exactly $1+4d+4q$;
no slack or additional tangent variables were added.
% END CRW18

% BEGIN CRW19
\paragraph{CRW19. The complete one-point ChainAbs argument.}
\label{claim:CRW19}
Let the source's actual patched data supply a
Noetherian ring $R_{\mathrm{pat}}$, relation ideal
$I_{\mathrm{pat}}$, annihilator $A_{\mathrm{pat}}$,
one prime $p_*$ containing $I_{\mathrm{pat}}$,
and the completion
\[
 C=\widehat{(R_{\mathrm{pat}}/I_{\mathrm{pat}})_{\bar p_*}}
\]
with its support ideal $S$.
Retain the exact fiber equivalence on every prime
$Q\subset C$ containing all source parameter images:
\[
 S\subset Q\quad\Longleftrightarrow\quad
 A_{\mathrm{pat}}\subset
       \operatorname{contraction}_{R_{\mathrm{pat}}}(Q).
                                                               \tag{CRW.29}
\]
Retain also that each minimal prime $Z$ over $S$
satisfies $\dim(C/Z)\ge1+D_0$.
Suppose this is the same $C$ and lifted ordinary
point to which CRW18 applies.

Assume a prime $p_0$ of $R_{\mathrm{pat}}/I_{\mathrm{pat}}$
lies below $\bar p_*$, contains all parameter images,
and has contraction to $R_{\mathrm{pat}}$ containing
$A_{\mathrm{pat}}$.
It lifts to a prime $Q_0$ of $C$: localize at
$\bar p_*$ and use faithful flatness of completion.
This theorem applies because the localized ring is
Noetherian local.  The contraction is exactly $p_0$.
By (CRW.29), $S\subset Q_0$.
There is a minimal prime $Z$ over $S$ below $Q_0$:
localize the quotient $C/S$ at $Q_0/S$ and take a
minimal prime, or use existence of minimal primes
in a Noetherian ring.
The dimension hypothesis and CRW18 give $Z=0$.
In particular $S=0$.

For every target prime $p$ below $\bar p_*$ containing
the parameters, lift it to a prime $Q$ of $C$ in
the same way.  Now $S=0\subset Q$, so (CRW.29)
gives $A_{\mathrm{pat}}\subset\operatorname{contraction}(p)$.
This is exactly the $h=0$ conclusion of the original
$\mathrm{chainAbs}$.  Its initial stable condition
and target stable-or-thick condition may be kept
without alteration.  The stronger local comparison
has made them unnecessary inside this proof.
There is no connecting prime, no step to another
completion, and no use of the source's $\mathrm{move}$ field.
The support ideal and its dimension and fiber data
have not been manufactured by the argument: their
production is the explicit retained patching input.
% END CRW19

% BEGIN CRW20
\paragraph{CRW20. The exact conditional installation in Core.}
\label{claim:CRW20}
We spell out the join to the original patching construction,
to distinguish the completed implication from its remaining
producers.  Supply its original $\Lambda:\mathrm{Lim}(T)$,
$W:\mathrm{RRObj}$ and $m:\mathrm{RMaps}(W)$ on the same
$\mathrm{TWAut}(u,J,\mathrm{Unit},D,T)$.
Keep the actual splitting algebra $\mathrm{Spo}$, the
surjection $\pi:W.R/I_{\mathrm{pat}}\to\mathrm{Spo}$,
the lifted point $f':\mathrm{Spo}\to A$, and the ordinary
point $P$ on $W.R/I_{\mathrm{pat}}$.
Write $\pi_R:W.R\to\mathrm{Spo}$ for its composite
with the quotient map; this is the original $m.\pi$.
They must satisfy literally
\[
 P.f=f'\circ\pi,\qquad f'|_B=f,
 \qquad\ker f'=p_{\mathrm{Spo}},
\]
and preserve the same $x,k$, ordered characters and auxiliary
roots.  This is the data constructed by the source's
$\mathrm{rrEx}$, $\mathrm{rmapsEx}$, $\mathrm{piKer}$,
$\mathrm{piSurj}$, $\mathrm{lift\_pt}$ and $\mathrm{ooPt}$.
It is retained here, not inferred from the record names.

The original S5 output on these modules supplies
the completion/support data of CRW19.
Its parameter count is
$|\operatorname{Fin}(q)\times\operatorname{Fin}(n)|=nq$;
the regular coefficient completion has dimension $1+D$.
Hence its dimension lower bound is
\[
 (1+D)+nq=1+D_0,                                  \tag{CRW.30}
\]
exactly the one in CRW18.
CRW12--CRW15 supply K1 for this same $\Lambda.h$,
using (CRW.2).  CRW2--CRW18 apply to this literal
lifted point and completion.  Thus CRW19 supplies the
annihilator containment on the single completed support.

For descent to the original ideal, retain the actual source
equivalence on primes $Q$ of $\mathrm{Spo}$,
\[
 A_{\mathrm{pat}}\subset\pi_R^{-1}(Q)
    \quad\Longleftrightarrow\quad J\subset Q\cap B,
                                                               \tag{CRW.31}
\]
and the prime-lifting property of the finite free splitting
algebra used by $\mathrm{chain\_lift}$.
Given a starting prime $q_1\subset\ker f$ containing
$J$, and any target prime $q\subset\ker f$, lift them
below $p_{\mathrm{Spo}}$, apply (CRW.31) to the starting
prime, then CRW19, then (CRW.31) to the target prime.
The result is $J\subset q$.
All parameter images vanish in these primes by the exact
$m.\pi_s$ identity.  This proves the one-point Core
implication from the explicitly enumerated patching data.
It uses neither a second splitting prime chosen for a
different point nor a separate ideal for the support step.
% END CRW20

% BEGIN CRW21
\paragraph{CRW21. Completed producers, remaining Raw cut and common witnesses.}
\label{claim:CRW21}
The new completed subpackets are the actual full-block
S6/URep (CRW2--CRW6); selected CohLim on a supplied original
limit (CRW12--CRW15); actual cotangent generators and
power-series presentation (CRW7--CRW18); and the
one-completion support implication (CRW19--CRW20).
In particular the full-block representation and its
bounded derivative are no longer additional open inputs
at this endpoint.  The relative cotangent calculation
does not add any hypothesis about $\mathrm{Sel1}=0$.

The remaining sufficient cut for the original selected
Raw instance is still substantive.  It must supply
the original pro-$\ell$ quotient and its descent of
the pair, residual data, point, B1/B2, same ideal and
TWAut; the coefficient normalization and added Unit
parameter of CoreBC; the actual limit S4; and the
compatible patched ring/maps/modules just enumerated,
including the S5 support dimension and fiber equivalence.
In S5 these latter fields come from eventual freeness
modulo every $(k,j)$, the two quotient comparison maps
with composites multiplication by the same $f$, one
uniform exponent for approximate annihilation of every
relation, and the regular coefficient completion.
None of them follows just from the local power-series
presentation.  The current packet does not assert
that those producers have been reconstructed.

When that cut is supplied, the selected Raw statement
has $h=0$ and one prime $p=\ker f$.
JSP10 supplies its original stable discriminant.
Taking the starting prime to be $p$ shows
\[
 J\subset p\quad\Longrightarrow\quad
   \forall q_0\text{ prime},\ q_0\subset p
             \Longrightarrow J\subset q_0.        \tag{CRW.32}
\]
This is precisely the NiceI consumer in
$\mathrm{PFromAutOf}$ needed for P1core.
There is no arbitrary-Kap, NiceE or P3core assertion.
The single-point residue-cardinality comparison is
reflexive, but the old coefficient map need not be
onto the point residue field: CoreBC's normalization
must still be supplied before CRW2 applies.

The ideal $J$ remains the kernel/annihilator fixed by
the same original ordinary Hecke system in JSP.
The ordinary/control/type-neat-level and TThetaEx
producer, the four further Hida endpoints on this $J$,
and the all-level HasAut on this identical $J$ remain
open.  CohP, LocCFT, the original N4 old-subgroup
condition and the complete arithmetic setting remain
the inherited hypotheses; the new local arguments
have not silently replaced them.
This installation must be made for both original
rank-four Links.  Fix the complex embedding $\iota$
first, then choose the one common coefficient threshold
for both Links and both properties on each ideal.
The point, extension, detector, $k$, $T$, TWPl, TWAut
and $J$ stay fixed through all uses above.
The constant in (CRW.2) lies inside those fixed data,
before $m$ and eventual $N$, and never replaces the
outer common threshold.

This advances the Count-II branch required by both
RH and unrestricted Dirichlet GRH.  KA5--KA6 still
remove only Case III at two-power-order targets,
including the principal target.  Counts I, II and V
remain required for RH; the extra Case-III producer
remains required for unrestricted Dirichlet GRH.
% END CRW21

\subsection{Published inputs and exact applicability}

The only non-elementary representation input here is
Chenevier's Cayley--Hamilton result in CRW3.  Its
primary text is \emph{The $p$-adic analytic space of
pseudocharacters of a profinite group and
pseudorepresentations over arbitrary rings},
Lemma 1.21, \S1.22 and Theorem 2.22(i), printed
pages 18--19 and 34 (PDF pages 18--19 and 34):
\url{https://maths.dur.ac.uk/events/Meetings/LMS/2011/GRAF11/papers/chenevier.pdf}.
The hypotheses and full-law scope are checked in CRW3;
CRW4--CRW6 prove the additional normalization and
closed-ideal properties independently.
The complete-local-to-henselian implication is
Lemma 10.153.9 of the Stacks Project
(\url{https://stacks.math.columbia.edu/tag/04GM});
the actual ring $R$ is complete and separated.

The following commutative-algebra applications are also
checked against primary statements, with separate
certificates in \texttt{published\_applications.json}.
The Stacks Project, Lemmas 10.31.1--10.31.2
(\url{https://stacks.math.columbia.edu/tag/00FM})
apply to the Noetherian DVRs and the finite-variable
power-series rings used in CRW2 and CRW17--CRW18.
Artin--Rees, Lemma 10.51.2
(\url{https://stacks.math.columbia.edu/tag/00IN}),
and Krull intersection, Lemma 10.51.4
(\url{https://stacks.math.columbia.edu/tag/00IP}),
apply in CRW5 to the finite modules $M\subset P_0$
over the actual Noetherian local $B$.
Lemma 10.97.3
(\url{https://stacks.math.columbia.edu/tag/00MC})
and Lemma 10.97.4
(\url{https://stacks.math.columbia.edu/tag/031C})
apply to the maximal-ideal completions of the
Noetherian localizations in CRW2, CRW11 and CRW19.
Lemma 15.44.5
(\url{https://stacks.math.columbia.edu/tag/0AP1})
applies to the DVR $O[[S]]_{(\varpi)}$ in CRW2.
Proposition 10.60.9
(\url{https://stacks.math.columbia.edu/tag/00KQ})
applies to the local Noetherian power-series ring
in CRW18, whose maximal ideal has $1+D_0$ generators.
These applications do not supply the missing patched
support ideal, its dimension lower bound, or the
all-level automorphic modules.
